\documentclass[10pt,a4paper]{amsart}
\usepackage[margin=19mm]{geometry}
\usepackage{amsmath,amssymb,mathtools}
\usepackage[scr=rsfs,cal=euler]{mathalfa}
\usepackage{xfrac}
\usepackage[unicode,hidelinks,bookmarks=false]{hyperref}
\newcommand{\ie}{i.e.\ }
\newcommand{\cf}{cf.\ }
\newcommand{\resp}{respectively\ }
\newcommand{\Subsection}{\subsection}
\providecommand{\nobreakdash}{\nobreak\hbox{-}}
\setlength{\emergencystretch}{2em}
\multlinegap0pt
\usepackage{amssymb}
\usepackage{enumerate}
\usepackage{xcolor}
\usepackage{tikz}
\usepackage{bbm}
\usepackage{multirow}
\usepackage[shortlabels]{enumitem}
\newtheorem{prop}{Proposition}
\newtheorem{lemma}{Lemma}
\newcommand{\A}{\mathcal{A}}
\newcommand{\ra}{\rightarrow}
\title{Symplectic Hodge Theory on Lie Algebroids}
\author{Shane Rankin}
\date{Source: arXiv:2411.06012v2 (CC BY 4.0)}
\begin{document}
\maketitle

\noindent\emph{Source and licence.} Symplectic Hodge Theory on Lie Algebroids. Version arXiv:2411.06012v2 (CC BY 4.0), distributed by arXiv under the Creative Commons Attribution 4.0 International licence, http://creativecommons.org/licenses/by/4.0/, as recorded in the arXiv licence metadata for this exact version and shown on the abstract page https://arxiv.org/abs/2411.06012v2. Full licence notice: \texttt{licenses/arxiv-2411.06012-CC-BY-4.0.txt}.\par\smallskip

\noindent\textit{Editorial scope.} Counted unit: the article's lemma AlgebroidCommutators (source0.tex lines 1144--1155). The article's complete original proof of it is delivered with it (source0.tex lines 1156--1203, one \texttt{proof} environment of 48 lines). No mathematics is written, repaired or re-derived: the statement and the proof are the article's own lines, in the article's own notation, and no retained line is substituted or re-flowed.  Retained uncounted as the source-local context the proof needs: lines 1109--1111; lines 1136--1138.  Nothing was omitted from any retained span, so the package carries no omission record.  No retained line contains a citation, so the wrapper carries no bibliography and no bibliography entry.  No retained line contains a figure environment, an image-inclusion command or drawing code, so no image is delivered and the package declares no asset dependency.  Editorial formatting only, in addition: an AMS \texttt{amsart} preamble that restates the article's own declarations and macros used by the retained lines, namely the environments \texttt{prop}, \texttt{lemma} on separate counters, together with the standard packages those lines need.  The retained environments print in this document as Proposition 1, Lemma 1 in order of appearance, whereas the article prints the same items with the numbering of its own full document.  The complete native source of the article as distributed in the arXiv e-print for this exact version is bundled unchanged as \texttt{sources/168248/source0.tex}.
\begin{prop}\label{CartanWedge2}
    Given $X,Y \in \Gamma(\wedge^2 \A)$, we have that $[\iota_X, [\iota_Y,d_\A]]=\iota_{[X,Y]}$ where $[\cdot,\cdot]$ is the Nijenhius-Schouten extension of the bracket on $\A$.
\end{prop}

\begin{prop}\label{PoissonBracketVanish}
    Let $(\A \ra M,\rho,\omega)$ be a symplectic Lie algebroid, and let $\pi \in \Gamma(\bigwedge^2\A)$ be the dual bivector to $\omega$ as above. Then $[\pi,\pi]=0$.
\end{prop}

\begin{lemma}\label{AlgebroidCommutators}
    Let $(\A \ra M,\rho,\omega)$ be a symplectic Lie algebroid. Then we have the following commutation relations:
    \begin{enumerate}
        \item $[L,d_\A]=0$.
        \item $[H,d_\A]=d_A$.
        \item $[H,d_\A^\star]=-d_\A^\star$.
        \item $[\Lambda , d_\A^\star]=0$.
        \item $[L,d_\A^\star]=d_\A$.
        \item $[d_\A,d_\A^\star]=0$.
        \item $[d_\A^\star,d_\A^\star]=0$.
    \end{enumerate}
\end{lemma}

 \begin{proof}
        We prove these in order for $\alpha \in \Omega_{\A}^k(M)$:
        \begin{enumerate}
            \item Using that $\omega$ is closed we have
            \begin{align*}
                [L,d_\A]\alpha &= \omega \wedge d_\A\alpha - d_\A(\omega \wedge \alpha) \\
                &= \omega \wedge d_\A \alpha - (d_\A\omega \wedge \alpha  +\omega \wedge d_\A \alpha) \\
                &= 0.
            \end{align*}
            \item Using only degree considerations we have
            \begin{align*}
               [H,d_\A]\alpha = Hd_\A \alpha  -d_\A H\alpha = [(k+1)-n]d_\A \alpha - (k-n)d_\A \alpha = d_\A \alpha.
            \end{align*}
            \item Again using only degree considerations we have that
            \begin{equation*}
                [H,d_\A^\star]=Hd_\A^\star \alpha - d_\A^\star H \alpha = [(k-1)-n]d_\A^\star \alpha - (k-n)d_\A^\star \alpha = -d_\A^\star \alpha.
            \end{equation*}
            \item Using Proposition \ref{CartanWedge2}, we have that
            \begin{equation*}
               [\Lambda,d_\A^\star] =[\iota_\pi,[\iota_\pi,d_\A]]=\iota_{[\pi,\pi]}.
            \end{equation*}
            To show this vanishes, we must have that $[\pi,\pi]=0$, but this follows from Proposition \ref{PoissonBracketVanish}.
            \item Using the fact that $\omega_\A$ is $d_\A$-closed, and that $[L,d_\A]=0$, we have
            \begin{align*}
                [L,d_\A^\star]\alpha &= Ld_\A^\star \alpha - d_\A^\star L\alpha \\
                &=L\Lambda d_\A \alpha - Ld_\A \Lambda \alpha - \Lambda(\omega_\A \wedge d_\A \alpha) + d \Lambda (\omega_\A \wedge \alpha) \\
                &=[L,\Lambda]d_\A \alpha - d[L,\Lambda]\alpha \\
                &= [H,d_\A]\alpha \\
                &=d_\A\alpha.
            \end{align*}
            
            \item The graded Jacobi identity tells us that
            \begin{equation*}
                -[d_\A,[\Lambda,d_\A]]+[\Lambda,[d_\A,d_\A]]+ [d_\A,[d_\A,\Lambda]]=0.
            \end{equation*}
            From this we must have that
            \begin{equation*}
                -2[d_\A,d_\A^\star]=0.
            \end{equation*}
            
            \item Using items 4 and 6, along with the graded Jacobi identity we have that
            \begin{align*}
                [d_\A^\star,d_\A^\star]&=[d_\A^\star,[\Lambda, d_\A]] \\
                &=[\Lambda, [d_\A,d_\A^\star]]+[d_\A, [d_\A^\star, \Lambda]] \\
                &=0.
            \end{align*}
        \end{enumerate}
    \end{proof}

\end{document}
